\section{Evaluation}
\label{sec:evaluation}

The evaluation asks two bounded questions: whether the production controller
preserves the required START--PATTERN--END ordering under randomized delays and
faults, and whether the same RTL is implementable through a practical FPGA
observability boundary. It does not measure link throughput, latency, BER, or
electrical behavior.

\subsection{Environment and Reproducibility}

Simulation used Questa Intel FPGA Starter Edition 2023.3. The committed
PowerShell and Questa scripts compile the RTL, UVM packages, and assertions,
then run fixed seed lists. The integrated testbench drives a 100-MHz simulation
clock, reduces the training sample count to eight accepted samples, and fixes
\texttt{ALLOW\_ABSTRACT\_DATATRAIN\_BYPASS} to zero. These changes shorten the
experiment without changing the message-order, accepted-sample, threshold, or
abort checks. The RTL default remains 4096 accepted samples.

FPGA mapping, fitting, and timing analysis used Quartus Prime Lite
23.1std.1, targeting a Cyclone 10 LP 10CL025YU256C8G. The SDC applies a
12.5-ns (80-MHz) clock constraint and an asynchronous-reset false path. It does
not assign package locations or external input/output delays; timing values are
therefore internal-clock evidence, not board timing closure.

\subsection{Production Integrated-Training Campaign}

Five seeds each execute 36 trials. Every seed first guarantees all nine
scenario classes and then adds seeded variations in response latency,
backpressure, sample gaps, corruption, threshold, and abort phase. Table
\ref{tab:integrated} reports the accepted per-seed summaries. A request need not
receive a response in malformed-response, retry, or abort cases, and END is
intentionally absent when the pattern does not pass. The resulting count
asymmetry is therefore expected, not packet loss.

\begin{table*}[t]
\caption{Production Integrated-Training Results With Bypass Disabled}
\label{tab:integrated}
\centering
\small
\setlength{\tabcolsep}{4.2pt}
\begin{tabular}{c r rr rr rr rr rrr r}
\toprule
Seed & Trials & \multicolumn{2}{c}{START} & \multicolumn{2}{c}{END} & Retry & Prot. & \multicolumn{2}{c}{Pattern} & \multicolumn{3}{c}{Abort} & Checks \\
 & & Req. & Resp. & Req. & Resp. & & err. & Pass & Fail & Reset & Fatal & T/O & \\
\midrule
1701 & 36 & 36 & 30 & 25 & 17 & 8 & 6  & 25 & 4 & 8 & 1 & 4 & 83 \\
1802 & 36 & 36 & 28 & 23 & 12 & 8 & 11 & 23 & 2 & 3 & 3 & 7 & 76 \\
1903 & 36 & 32 & 28 & 21 & 14 & 5 & 6  & 23 & 4 & 3 & 6 & 7 & 78 \\
2004 & 36 & 34 & 26 & 24 & 18 & 8 & 6  & 23 & 6 & 2 & 3 & 7 & 80 \\
2105 & 36 & 35 & 28 & 27 & 17 & 11& 6  & 27 & 4 & 4 & 5 & 4 & 81 \\
\midrule
Total & 180 & 173 & 140 & 120 & 78 & 40 & 35 & 121 & 20 & 20 & 18 & 29 & 398 \\
\bottomrule
\end{tabular}
\end{table*}

The scenario-hit totals were 18 nominal passes, 20 failed-pattern retry/pass
cases, 17 START retries, 16 malformed START responses, 23 END retries, 19
malformed END responses, 20 reset aborts, 18 fatal aborts, and 29 state-timeout
aborts. Thus all planned classes were observed, while the predictor performed
398 independent order/value/cause checks. The accepted logs report zero UVM
errors and zero UVM fatals, and the regression summary reports no assertion
failures. These observations establish consistency for the committed stimuli;
they are not an exhaustive proof.

\subsection{Preserved Regression Gate}

The integrated campaign was added without discarding earlier evidence. The
release gate also ran nine deterministic UVM tests, 200 sideband trials over
seeds 101--505, 160 LFSR-training trials over seeds 701--1105, and 180 recovery
trials over seeds 1201--1605. Module-directed suites separately exercise the
controller, sideband sequencer, LFSR engine, error manager, and wrapper CSR.
All accepted runs completed without UVM errors/fatals or reported assertion
failures. An ad-hoc ignored debug log is not part of this evidence set.

\subsection{FPGA Mapping, Fit, and Timing}

Table~\ref{tab:fpga} preserves both the positive and negative implementation
results. The wide controller top maps to 763 logic
elements and 505 registers, but demands all 151 device pins; its fitter run
fails. Timing numbers emitted alongside that failed fit are not used as
successful-fit evidence.

The CSR wrapper fits successfully using 824 of 24,624 logic elements (3\%), 505
registers, and 119 of 151 pins (79\%), with no virtual pins. Across the committed
timing corners, the worst setup and hold slacks are +1.013 and +0.179 ns,
respectively, and reported negative slack is zero. The wrapper adds a small
amount of decode logic while reducing external pin demand by 32, demonstrating
why serializing diagnostics is necessary for this device. Because external I/O
constraints and physical pin locations are absent, the result supports
internal synthesizability and fit only \cite{quartusUG}.

\begin{table}[t]
\caption{Cyclone 10 LP Implementation Summary}
\label{tab:fpga}
\centering
\small
\begin{tabular}{lrrr}
\toprule
Metric & Core map & Core fit & Wrapper fit \\
\midrule
Status & success & failure & success \\
Logic elements & 763 & 764 & 824 \\
Registers & 505 & 505 & 505 \\
Pins & 151 & 151 & 119 \\
Setup slack (ns) & -- & -- & +1.013 \\
Hold slack (ns) & -- & -- & +0.179 \\
\bottomrule
\end{tabular}
\end{table}

Fig.~\ref{fig:fpga-evidence} retains the successful and failed implementation
outcomes in one qualified data graphic. Reporting the failed wide-top fit is
important: omitting it would hide the pin-cost reason for the CSR wrapper.

\begin{figure*}[t]
\centering
\includegraphics[width=0.92\textwidth]{fig08_fpga_results.pdf}
\caption{Quartus implementation evidence for the same controller under two
observability boundaries. Timing is reported only for the successfully fitted
wrapper and is qualified as internal-clock evidence.}
\label{fig:fpga-evidence}
\end{figure*}

The complete release gate is summarized by Table~\ref{tab:integrated}, the
implementation results in Table~\ref{tab:fpga}, and a
\href{https://github.com/Varadain/ucie-ltsm/blob/feature/v1.0-integrated-training/journal/figures/svg/fig12_regression_evidence.svg}{supplementary release-gate graphic}.
Trial counts are not combined into a coverage percentage because the campaigns
exercise different interfaces and scenario spaces. The supplementary graphic
records which directed, deterministic, randomized, predictor, assertion, and
implementation evidence was preserved together without duplicating those data
as a full-width manuscript float.

\subsection{Threats to Validity}

The pseudo-random campaign uses explicit seeded \texttt{\$urandom\_range}
choices and sampled counters because the Starter license does not support the
native class-randomization and UCDB coverage flow required by the original
plan. No native functional-coverage percentage or code coverage is claimed.
The independent predictor reduces common-mode risk, but design and verification
were developed in the same project. Reduced training counts accelerate
simulation; they do not test full-duration drift. Finally, FPGA synthesis
cannot substitute for a complete two-die model, board experiment, ASIC PPA, or
electrical characterization.
